% Textbook Section 2.3, pp. 145--150.
\section[Invertibility]{Characterizations of Invertible Matrices}
\sectioncard{2.3}{Characterizations of Invertible Matrices}

\begin{frame}{Review Case: Is $A$ Invertible?}
Textbook Example 1 asks whether
\[
A=\begin{bmatrix}
1&0&2\\
3&1&2\\
5&1&9
\end{bmatrix}
\]
is invertible.
\pause
\begin{alertblock}{Review question}
What facts could establish invertibility without explicitly calculating $A^{-1}$?
\end{alertblock}
\end{frame}

\begin{frame}[shrink=11]{The Invertible Matrix Theorem}
\theorembox{Theorem 8}{Let $A$ be a square $n\times n$ matrix. The following statements are equivalent.}
\vspace{-.4ex}
\begin{columns}[T]
\column{.49\textwidth}
\begin{enumerate}[a.]
  \item $A$ is invertible.
  \item $A$ is row equivalent to $I_n$.
  \item $A$ has $n$ pivot positions.
  \item $A\vect x=\vect0$ has only the trivial solution.
  \item The columns of $A$ are linearly independent.
  \item $\vect x\mapsto A\vect x$ is one-to-one.
\end{enumerate}
\column{.49\textwidth}
\begin{enumerate}[a.]
  \setcounter{enumi}{6}
  \item $A\vect x=\vect b$ has a solution for every $\vect b\in\R^n$.
  \item The columns of $A$ span $\R^n$.
  \item $\vect x\mapsto A\vect x$ maps $\R^n$ onto $\R^n$.
  \item Some $C$ satisfies $CA=I$.
  \item Some $D$ satisfies $AD=I$.
  \item $A^T$ is invertible.
\end{enumerate}
\end{columns}
\pause

\vspace{0.1in}

\centering
\textbf{For a fixed square matrix, these statements are either all true or all false.}
\end{frame}

\begin{frame}{Textbook Example 1: Echelon Form}
\small
\[
\begin{bmatrix}
1&0&2\\3&1&2\\5&1&9
\end{bmatrix}
\xrightarrow[ R_3\leftarrow R_3-5R_1]{R_2\leftarrow R_2-3R_1}
\begin{bmatrix}
1&0&2\\0&1&-4\\0&1&-1
\end{bmatrix}
\]
\pause
\[
\xrightarrow{R_3\leftarrow R_3-R_2}
\begin{bmatrix}
1&0&2\\0&1&-4\\0&0&3
\end{bmatrix}.
\]
\pause
The echelon form has three pivot positions. By statement (c) of the Invertible Matrix Theorem, $A$ is invertible.
\end{frame}

\begin{frame}[shrink=6]{What Else Does the Echelon Form Tell Us?}
  \begin{center}
    \[
\begin{bmatrix}
1&0&2\\0&1&-4\\0&0&3
\end{bmatrix}.
\]
  \end{center}
\begin{columns}[T]
\column{.48\textwidth}
\textbf{Equations and columns}
\begin{itemize}
  \item Does $A\vect x=\vect0$ have a nontrivial solution? \pause No.
  \item Are the columns linearly independent? \pause Yes.
  \item Do the columns span $\R^3$? \pause Yes.
  \item Does $A\vect x=\vect b$ have a unique solution for every $\vect b$? \pause Yes.
\end{itemize}
\column{.48\textwidth}
\textbf{Transformations and inverses}
\begin{itemize}
  \item Is $\vect x\mapsto A\vect x$ one-to-one? \pause Yes.
  \item Is it onto $\R^3$? \pause Yes.
  \item Do left and right inverses exist? \pause Yes, and both equal $A^{-1}$.
  \item Is the inverse transformation linear? \pause Yes: $S(\vect x)=A^{-1}\vect x$.
\end{itemize}
\end{columns}
\end{frame}

\begin{frame}{Invertible Linear Transformations}
\theorembox{Theorem 9}{Let $T:\R^n\to\R^n$ have standard matrix $A$. Then $T$ is invertible if and only if $A$ is invertible. In that case,
\[
T^{-1}(\vect x)=A^{-1}\vect x.
\]}
\pause
\begin{block}{Textbook Example 2}
If $T:\R^n\to\R^n$ is one-to-one, the columns of $A$ are linearly independent. The Invertible Matrix Theorem then shows that $T$ is also onto and invertible.
\end{block}
\end{frame}

\begin{frame}{Transpose and One-Sided Inverses}
The Invertible Matrix Theorem includes
\[
A\text{ is invertible}\quad\Longleftrightarrow\quad A^T\text{ is invertible}.
\]
\pause
It also converts a one-sided inverse into a two-sided inverse. If square matrices satisfy
\[
AB=I,
\]
then $A$ and $B$ are both invertible, with
\[
B=A^{-1},\qquad A=B^{-1}.
\]
\end{frame}

% \begin{frame}[shrink=8]{Why the Statements Are Equivalent}
% The proof begins with a circle of implications:
% \[
% \boxed{(a)}\Rightarrow\boxed{(j)}\Rightarrow\boxed{(d)}
% \Rightarrow\boxed{(c)}\Rightarrow\boxed{(b)}\Rightarrow\boxed{(a)}.
% \]
% \pause
% The remaining statements attach to this circle through familiar equivalences:
% \[
% \boxed{(g)}\Longleftrightarrow\boxed{(h)}\Longleftrightarrow\boxed{(i)},
% \qquad
% \boxed{(d)}\Longleftrightarrow\boxed{(e)}\Longleftrightarrow\boxed{(f)}.
% \]
% \[
% \boxed{(a)}\Longleftrightarrow\boxed{(l)},
% \qquad
% \boxed{(a)}\Rightarrow\boxed{(k)}\Rightarrow\boxed{(g)}\Rightarrow\boxed{(a)}.
% \]
% \pause
% Because every statement connects back to the same circle, proving any one of them proves all twelve.
% \end{frame}

\begin{frame}[shrink=4]{Singular Matrices and Numerical Caution}
\begin{columns}[T]
\column{.47\textwidth}
\textbf{Textbook Practice Problem 1}
\[
A=\begin{bmatrix}
2&3&4\\2&3&4\\2&3&4
\end{bmatrix}.
\]
The columns are linearly dependent, so $A$ is singular. Equivalently, $A$ has fewer than three pivots and is not row equivalent to $I_3$.

\medskip
The theorem applies only to square matrices.
\column{.49\textwidth}
\textbf{Nearly singular matrices}

An invertible matrix can become singular after a very small change in its entries. Roundoff may then make a singular matrix appear invertible, or an invertible matrix appear singular.

\medskip
The condition number measures this sensitivity:
\begin{itemize}
  \item $\operatorname{cond}(I)=1$;
  \item a singular matrix has infinite condition number;
  \item a large condition number warns of possible numerical error.
\end{itemize}
\end{columns}
\end{frame}
